\section{Related Work}
\label{sec:related}

\subsection{3D Representations for VLA}
\label{sec:rw_spatial}

Vision-Language-Action models trained on 2D image observations struggle with tasks that require precise spatial reasoning, and a growing body of work seeks to close this gap by injecting 3D geometric structure into VLA architectures. These methods differ in the type of spatial representation they adopt, and we organize them accordingly.

\paragraph{Depth-based approaches.}
The most direct strategy augments VLA inputs or supervision with depth information. DepthVLA~\cite{an2025depthvla} conditions action prediction on monocular depth maps fused with RGB features, while QDepth-VLA~\cite{qdepthvla2025} quantizes depth predictions into discrete tokens that serve as auxiliary supervision during training. AugVLA-3D~\cite{augvla3d2026} augments vision encoder features with depth-driven spatial embeddings, and 3D-CAVLA~\cite{3dcavla2025} uses depth and 3D context to generalize across unseen tasks. These depth-based methods consistently improve spatial precision over RGB-only baselines, with gains concentrated on tasks that require reaching or placing objects at precise locations. Yet depth maps remain viewpoint-dependent and encode only surface geometry; they do not explicitly represent free-space boundaries or the volumetric extent of obstacles, information that is critical for collision avoidance in cluttered workspaces.

\paragraph{Point cloud and geometric representations.}
A second family encodes the workspace as explicit 3D geometry. PointVLA~\cite{pointvla2025} injects point cloud tokens into the VLA backbone to provide direct geometric grounding for action prediction, and GeoVLA~\cite{geovla2025} encodes geometric features from multi-view reconstructions into a unified spatial representation. GeoAware-VLA~\cite{geoaware2025} uses implicit geometry-aware representations without requiring point cloud inputs at inference time, offering a middle ground between explicit and implicit approaches. These methods capture richer 3D structure than depth maps but typically require multi-view cameras or known camera extrinsics, and the computational cost of processing raw point clouds at each timestep can be substantial for real-time control.

\paragraph{Implicit spatial alignment.}
Rather than providing explicit 3D inputs, several methods align the VLA's internal representations with spatial structure through auxiliary objectives or architectural inductive biases. SpatialVLA~\cite{wu2025spatialvla} introduces adaptive spatial grounding tokens that encode egocentric spatial relations across embodiments. Spatial Forcing~\cite{li2025spatialforcing} aligns vision-language features with spatial representations through an implicit contrastive objective and achieves strong generalization (ICLR~2026). Evo-0~\cite{evo02025} builds implicit spatial understanding via a progressive training curriculum. This category avoids runtime overhead from explicit 3D processing, at the cost of less interpretable spatial representations.

\paragraph{Structured spatial tokens and anchors.}
GST-VLA~\cite{chen2026gstvla} represents scenes through Gaussian spatial tokens that encode local 3D geometry in a differentiable form. StemVLA~\cite{stemvla2026} predicts future 3D spatial geometry as an auxiliary task, combining historical 4D representations with forward-looking spatial reasoning. CorridorVLA~\cite{corridorvla2026} constrains generative action heads through sparse spatial anchors that define safe motion corridors.

\paragraph{Spatiotemporal and grounded representations.}
ST-VLA~\cite{stvla2026} extends spatial reasoning to the temporal dimension with 4D spatiotemporal features for general manipulation. SG-VLA~\cite{sgvla2026} learns spatially-grounded representations for mobile manipulation by binding visual features to metric locations. SA-VLA~\cite{savla2026} combines spatially-aware representations with flow-matching action generation for reinforcement learning.

Despite the diversity of spatial representations explored in these works, they share a common evaluation focus: task success rate and generalization to new objects, scenes, or embodiments. None of these methods evaluate collision safety as an outcome metric, and the question of whether improved spatial representations translate to fewer unintended collisions remains unanswered.

\subsection{Safe Robot Manipulation with VLAs}
\label{sec:rw_safe}

A parallel line of research addresses the safety of VLA-based manipulation policies, focusing on collision avoidance and constraint satisfaction during task execution. As VLA models are deployed in unstructured environments with humans and fragile objects, safety failures carry real consequences: a single unintended collision can damage objects, injure people, or erode user trust. These methods differ in how they integrate safety mechanisms into the VLA pipeline.

\paragraph{Plug-in barrier functions.}
VLSA~\cite{hu2025aegis} (also known as AEGIS) introduces a plug-and-play safety constraint layer that constructs a Control Barrier Function (CBF) from a world model's state estimates and filters the VLA's proposed actions to avoid collisions. This approach treats the VLA as a black box and applies safety constraints at the action level, achieving a 59.16\% improvement in Collision Avoidance Rate on the SafeLIBERO benchmark. The barrier function operates on Euclidean distances between object positions, which requires accurate state estimation and limits the safety guarantee to scenarios with known object geometries.

\paragraph{Constrained reinforcement learning.}
SafeVLA~\cite{zhan2025safevla} formulates VLA safety as a constrained optimization problem and trains the policy with reward shaping and constraint penalties that discourage unsafe behaviors. Selected as a NeurIPS~2025 Spotlight, SafeVLA demonstrates that RL-based alignment can produce safer policies without modifying the VLA architecture. The approach requires online interaction data and task-specific reward engineering for each safety specification, which limits its scalability to new domains where collision definitions vary.

\paragraph{Risk estimation and multi-modal perception.}
CoFreeVLA~\cite{chen2026cofreevla} takes a different approach by semi-internalizing collision risk: a learned risk estimation module runs alongside the VLA policy and predicts collision probability for candidate actions so that the model can select safer alternatives without an external safety filter. This represents a step toward safety-aware perception, though the risk module operates on semantic features rather than explicit spatial representations. Safe-Night VLA~\cite{safenightvla2026} extends VLA safety to low-visibility conditions by fusing thermal imaging with RGB observations and applies CBF-based filtering to address the perceptual limitations of RGB-only policies in dark or occluded environments. A recent survey~\cite{vlasafetysurvey2026} provides a systematic taxonomy of threats, evaluation protocols, and defense mechanisms for VLA safety, and identifies the lack of spatial grounding in current safety methods as an open challenge.

These safety-focused methods operate on action-level constraints, reward signals, or risk scores computed from semantic features or explicit state estimates. None of them exploit the 3D spatial representations developed in Section~\ref{sec:rw_spatial}. The spatial representation backbone is treated as a fixed feature extractor, and the relationship between representation quality and collision outcomes is neither studied nor optimized.

\subsection{Occupancy Representations and Barrier Functions in Robotics}
\label{sec:rw_occ}

Occupancy-based representations and barrier functions have been studied independently in robotics, and recent work provides convergent evidence that occupancy-informed barriers outperform Euclidean-distance alternatives. OccVLA~\cite{liu2025occvla} applies volumetric occupancy prediction as auxiliary supervision for a VLA in autonomous driving; occupancy-grounded features improve spatial understanding at zero inference cost. OGM-CBF~\cite{ji2024ogmcbf} constructs Control Barrier Functions from occupancy grid maps via signed distance fields for mobile robot navigation, and its results confirm that grid-based spatial representations yield more conservative and reliable safety constraints than point-based distance metrics. In manipulation, Learning to Nudge~\cite{kwon2026learningnudge} trains composable Differentiable Control Barrier Functions (DCBF) that reason about individual objects and enable safe interaction in dense clutter. Neural CS-BF~\cite{neuralcsbf2025} learns barrier functions directly in configuration space for manipulation planning and eliminates the need for explicit workspace geometry.

These works establish that occupancy and barrier function methods each improve safety in their respective domains. OccVLA validates occupancy-based auxiliary supervision for VLAs but targets autonomous driving rather than manipulation safety. OGM-CBF and Learning to Nudge demonstrate effective barrier function designs but rely on hand-crafted geometric inputs or object-centric decompositions rather than learned VLA features. No existing work combines occupancy-grounded representation learning with safety-aware VLA training for manipulation. The interaction between spatial representation choices (depth, point cloud, occupancy) and collision safety outcomes has not been studied, leaving a gap that this paper addresses through controlled ablation of auxiliary spatial supervision and its effect on both task success and collision avoidance.
